\chapter{Implementación}

\paragraph*{}
Una vez definidos el análisis y el diseño de la ampliación, la implementación se centró en llevar esas decisiones al dashboard de forma práctica. En este capítulo no se busca describir cada función del código ni enumerar todas las variables utilizadas, porque ese nivel de detalle haría la lectura más pesada de lo necesario. La intención es explicar, con una visión más general, cómo se ha construido la ampliación, qué librerías han tenido más peso y de qué manera cada objetivo del proyecto se ha trasladado a la aplicación final.

\paragraph*{}
La base del sistema sigue siendo una aplicación desarrollada en Python con Dash. Sobre ella se han añadido nuevos modelos de supervivencia, nuevas visualizaciones, generación de informes, interpretación mediante inteligencia artificial y soporte multilingüe. La verdad es que la implementación se ha planteado como una ampliación progresiva del dashboard existente, conservando su funcionamiento principal y haciendo crecer sus posibilidades de análisis.

\section{Base tecnológica utilizada}

\paragraph*{}
El desarrollo se apoya en un conjunto de librerías que cubren tres necesidades principales: tratamiento de datos, análisis estadístico y presentación de resultados. Algunas ya formaban parte del proyecto original, como Dash, Pandas, NumPy, Plotly o lifelines. Otras se han incorporado o reforzado para poder cumplir los nuevos objetivos, especialmente en lo relativo a Random Survival Forest, exportación de informes e interpretación local mediante inteligencia artificial.

\begin{table}[H]
\centering
\small
\begin{tabularx}{\textwidth}{|p{0.28\textwidth}|X|}
\hline
\textbf{Tecnología o librería} & \textbf{Uso principal dentro de la implementación} \\ \hline
\texttt{Dash} & Construcción de la interfaz web, navegación entre páginas, actualización dinámica de resultados y conexión entre las acciones del usuario y la lógica interna de la aplicación. \\ \hline
\texttt{Pandas} y \texttt{NumPy} & Carga, limpieza y transformación del dataset, además de preparación de los datos que posteriormente se entregan a los modelos de supervivencia. \\ \hline
\texttt{lifelines} & Ajuste de modelos clásicos de análisis de supervivencia, como Kaplan-Meier, Cox, Exponencial y Weibull \cite{davidsonpilon2019lifelines}. \\ \hline
\texttt{scikit-survival} & Implementación de Random Survival Forest y cálculo de métricas asociadas al rendimiento del modelo \cite{polsterl2020scikit}. \\ \hline
\texttt{Plotly} & Generación de gráficas interactivas para el dashboard y reutilización de esas mismas visualizaciones en la exportación de informes. \\ \hline
\texttt{ReportLab} y \texttt{Kaleido} & Creación de documentos PDF y conversión de gráficas de Plotly en imágenes estáticas insertables en los informes. \\ \hline
\texttt{requests}, \texttt{llama.cpp} y Qwen & Comunicación con un modelo de lenguaje ejecutado en local para generar explicaciones textuales de los resultados \cite{llamacpp_repo,llamacpp_server,qwen25_technical_report,qwen25_huggingface_gguf}. \\ \hline
\end{tabularx}
\caption{Tecnologías principales utilizadas en la implementación}
\label{tab:tecnologias_implementacion}
\normalsize
\end{table}

\paragraph*{}
Con este conjunto de herramientas se cubre el flujo completo de trabajo: preparación del dataset, ajuste de modelos, visualización de resultados, generación de informes e interpretación textual. Cada librería aporta una pieza concreta, lo que evita concentrar toda la lógica en una única parte de la aplicación.

\section{Estructura del código fuente}

\paragraph*{}
El código fuente se organiza en varios módulos, cada uno con una responsabilidad principal dentro del dashboard. No se detalla aquí el contenido interno de cada función, porque eso haría el capítulo demasiado técnico, pero sí conviene mostrar cómo se reparte la aplicación. Esta estructura ayuda a entender que la ampliación no se ha construido como un bloque aislado, sino como un conjunto de piezas conectadas entre sí.

\begin{table}[H]
\centering
\small
\begin{tabularx}{\textwidth}{|p{0.34\textwidth}|X|}
\hline
\textbf{Archivo o grupo de archivos} & \textbf{Responsabilidad general} \\ \hline
\texttt{cargaDataset.py} y \texttt{layout.py} & Definen la aplicación Dash, la navegación principal y la estructura visual de las páginas del dashboard. \\ \hline
\texttt{analysis\_callbacks.py} & Conecta la interfaz con los análisis estadísticos, actualizando tablas, gráficas y mensajes según las acciones del usuario. \\ \hline
\texttt{preprocesamiento.py} & Prepara el dataset antes de ejecutar los modelos, realizando tareas de carga, limpieza, validación y transformación de datos. \\ \hline
\texttt{kaplan\_meier.py}, \texttt{cox\_regression.py} y \texttt{log\_rank\_test.py} & Contienen las técnicas de supervivencia principales y las ampliaciones asociadas a los análisis clásicos del sistema. \\ \hline
\texttt{exponential.py}, \texttt{weibull.py} y \texttt{rsf.py} & Agrupan los nuevos modelos incorporados en la ampliación: Exponencial, Weibull y Random Survival Forest. \\ \hline
\texttt{survival\_plots.py} & Centraliza parte de las visualizaciones estadísticas generadas con Plotly. \\ \hline
\texttt{ollama\_AI.py} & Gestiona la comunicación con el modelo local de inteligencia artificial y la generación de interpretaciones textuales. \\ \hline
\texttt{pdf\_callbacks.py} y \texttt{pdf\_exporter.py} & Permiten seleccionar resultados desde la interfaz y construir informes PDF con tablas, gráficas e interpretaciones. \\ \hline
\texttt{translations.py} y \texttt{config.py} & Recogen textos traducibles, parámetros generales, rutas y configuración necesaria para el funcionamiento del sistema. \\ \hline
\end{tabularx}
\caption{Estructura general del código fuente de la aplicación}
\label{tab:estructura_codigo}
\normalsize
\end{table}

\paragraph*{}
Esta distribución facilita localizar cada parte del sistema. Si en el futuro se quisiera añadir una nueva técnica o modificar una salida concreta, el cambio podría abordarse sobre un módulo concreto sin rehacer toda la aplicación.

\section{Implementación por objetivos}

\paragraph*{}
La ampliación se organizó alrededor de los objetivos planteados al inicio del trabajo. Cada mejora se integró en el punto donde tenía más sentido: modelos, visualización, exportación, traducción o interpretación de resultados. Así, las nuevas funcionalidades se perciben como parte natural del dashboard y no como añadidos aislados.

\begin{table}[H]
\centering
\small
\begin{tabularx}{\textwidth}{|p{0.27\textwidth}|X|p{0.23\textwidth}|}
\hline
\textbf{Objetivo de la ampliación} & \textbf{Implementación realizada} & \textbf{Librerías destacadas} \\ \hline
Incorporar nuevos modelos de supervivencia & Se añadieron modelos paramétricos, como Exponencial y Weibull, y un modelo de aprendizaje automático basado en Random Survival Forest. & \texttt{lifelines}, \texttt{scikit-survival} \\ \hline
Ampliar el análisis de variables & Se adaptó el preprocesamiento para trabajar con nuevas variables del dataset y preparar entradas válidas para los distintos modelos. & \texttt{Pandas}, \texttt{NumPy} \\ \hline
Mejorar la visualización de resultados & Se incorporaron nuevas curvas, comparaciones entre técnicas, gráficos de importancia y representaciones asociadas a Cox y Log-Rank. & \texttt{Plotly}, \texttt{Dash} \\ \hline
Exportar informes & Se desarrolló la generación de informes PDF con tablas, gráficas e interpretaciones seleccionadas desde la interfaz. & \texttt{ReportLab}, \texttt{Kaleido} \\ \hline
Añadir interpretación automática & Se integró un modelo local de lenguaje para explicar los resultados estadísticos de manera textual y controlada. & \texttt{requests}, \texttt{llama.cpp}, Qwen \\ \hline
Internacionalizar la aplicación & Se centralizaron los textos de la interfaz y de parte de los resultados para poder alternar entre castellano e inglés. & \texttt{Dash}, diccionarios Python \\ \hline
\end{tabularx}
\caption{Relación entre objetivos de ampliación, implementación y librerías utilizadas}
\label{tab:implementacion_objetivos}
\normalsize
\end{table}

\subsection{Nuevos modelos de supervivencia}

\paragraph*{}
Uno de los cambios principales ha sido la incorporación de nuevas técnicas de análisis de supervivencia. Los modelos Exponencial y Weibull se han implementado como alternativas paramétricas que permiten estudiar el comportamiento del abandono académico desde una perspectiva más ajustada a una forma matemática concreta. El modelo Exponencial resulta útil cuando se asume un riesgo constante en el tiempo, mientras que Weibull ofrece más flexibilidad al permitir que el riesgo aumente o disminuya.

\paragraph*{}
También se añadió Random Survival Forest, que aporta una visión diferente al no depender de una forma paramétrica cerrada. Este modelo permite trabajar con varias características del estudiante y estimar patrones de supervivencia desde un enfoque de aprendizaje automático. Para el usuario, la diferencia se traduce en una nueva página de análisis con métricas, curvas representativas y una lectura más predictiva del abandono.

\subsection{Preprocesamiento y variables de análisis}

\paragraph*{}
El preprocesamiento se ajustó para que el dataset quedara preparado antes de ejecutar cualquier técnica. En esta fase se limpian nombres de columnas, se validan campos necesarios, se eliminan registros problemáticos y se construye una representación más estable de cada estudiante. No se trata solo de dejar los datos ``bonitos'', sino de evitar que un modelo reciba información duplicada, incompleta o incoherente.

\paragraph*{}
Además, la ampliación permite trabajar con variables de distinto tipo. Algunas son binarias, otras categóricas y otras numéricas. Por eso, la implementación prepara los datos de manera diferente según la técnica que se vaya a ejecutar. Por ejemplo, ciertos análisis necesitan comparar grupos, mientras que otros pueden utilizar directamente una matriz de características más amplia.

\subsection{Visualización y comparación de resultados}

\paragraph*{}
Las visualizaciones se han construido con Plotly para mantener la interactividad propia del dashboard. Esta decisión es importante porque las gráficas no funcionan como simples imágenes pegadas en la pantalla, sino como objetos dinámicos que Dash puede actualizar cuando el usuario selecciona una técnica, cambia una variable o solicita una comparación.

\paragraph*{}
En esta parte se incorporan curvas de supervivencia, comparaciones entre modelos, gráficos de importancia de variables y representaciones específicas para Cox y Log-Rank. Además, la comparación de técnicas se ha integrado como una opción accesible desde la zona de análisis de supervivencia, de modo que el usuario puede revisar primero los modelos por separado y después consultar una visión conjunta.

\subsection{Interpretación mediante inteligencia artificial}

\paragraph*{}
La interpretación automática se ha planteado como una ayuda para leer los resultados, no como una sustitución del análisis estadístico. El sistema calcula primero las métricas y tablas mediante las librerías correspondientes, y después envía un resumen controlado al modelo de lenguaje local para obtener una explicación en prosa.

\paragraph*{}
La ejecución local se realiza mediante \texttt{llama.cpp} y un modelo Qwen en formato cuantizado. Esto permite generar textos sin depender de un servicio externo, algo especialmente interesante cuando se trabaja con resultados académicos o datos que conviene mantener dentro del entorno de trabajo. El resultado es una explicación más cercana para el usuario, útil sobre todo cuando las métricas pueden resultar algo frías o difíciles de interpretar a primera vista.

\subsection{Exportación de informes}

\paragraph*{}
La exportación a PDF se implementó para que los resultados del dashboard pudieran conservarse fuera de la aplicación. El informe recoge información del análisis seleccionado, tablas resumen, gráficas y, si el usuario lo desea, la interpretación generada mediante inteligencia artificial. De esta forma, el dashboard no solo sirve para explorar datos en pantalla, sino también para producir un documento que puede revisarse o compartirse posteriormente.

\paragraph*{}
En este proceso, Plotly y Kaleido tienen un papel importante. Las gráficas que se ven en el dashboard se convierten en imágenes estáticas para poder insertarlas en el PDF, mientras que ReportLab se encarga de construir la estructura del documento. Así se mantiene una continuidad clara entre lo que el usuario analiza en la interfaz y lo que aparece finalmente en el informe exportado.

\subsection{Internacionalización del dash}

\paragraph*{}
El soporte multilingüe se añadió centralizando los textos de la aplicación en un sistema de traducciones. Esto permite cambiar entre castellano e inglés sin duplicar las páginas ni crear versiones separadas de la interfaz. La implementación se apoya en diccionarios de Python y en el estado de idioma gestionado desde Dash.

\paragraph*{}
Aunque pueda parecer una mejora secundaria, la internacionalización afecta a muchas partes del dashboard: menús, botones, títulos, mensajes, tablas e incluso textos generados para algunos resultados. Por eso se implementó como una capa transversal, pensada para acompañar a las distintas vistas de la aplicación.

\section{Integración general en el dashboard}

\paragraph*{}
La integración práctica se realiza mediante callbacks de Dash. El usuario interactúa con botones, selectores o pestañas; la aplicación recupera los datos necesarios, ejecuta el análisis correspondiente y actualiza la vista con nuevas tablas, gráficas o mensajes. Este mecanismo común permite que modelos distintos se utilicen de forma parecida desde la interfaz, aunque internamente cada técnica requiera un tratamiento propio.

\section{Resultado de la implementación}

\paragraph*{}
En conjunto, la implementación mantiene la base del sistema original y la amplía con mayor capacidad analítica, informes descargables, soporte multilingüe e interpretación asistida. El resultado es un dashboard más completo para estudiar el abandono académico universitario, sin perder la continuidad con la aplicación de partida.
